%======================
% Chapter 2 : Literature Review
%======================

\chapter{\scshape Literature Review}

Fall detection using wearable inertial sensors has been an active research area for over two decades. The literature can be broadly categorized by detection approach: threshold-based methods, machine learning classifiers, and more recent deep learning architectures.

\textbf{Threshold-based methods} compute a scalar magnitude from accelerometer axes --- typically the resultant $|a| = \sqrt{a_x^2 + a_y^2 + a_z^2}$ and trigger a fall event when a spike followed by a quiet period is detected. Kangas et al.\ \cite{kangas2008comparison} compared multiple threshold algorithms on hip and wrist IMU data and reported sensitivity values of 73--95\% depending on placement, but with false-positive rates exceeding 30\% on vigorous ADLs. Bourke and Lyons \cite{bourke2008threshold} proposed a dual-axis threshold method achieving 96\% sensitivity, though specificity on sit-to-stand transitions remained problematic. The fundamental limitation of threshold-only approaches is that any rapid motion --- vigorous exercise, dropping an object, stumbling without falling --- can trigger the threshold, requiring either a very high threshold (missing genuine falls) or a low threshold (excessive false positives).

\textbf{Machine learning classifiers} applied to hand-crafted time-domain and frequency-domain features from IMU windows have addressed some of these limitations. Sucerquia et al.\ published the SisFall dataset \cite{sucerquia2017sisfall}, a widely used benchmark containing 15 fall types and 19 ADL types from 38 subjects, enabling standardized evaluation of feature-based classifiers. Support vector machines, decision trees, and random forests trained on SisFall features have achieved sensitivity above 90\%, but these approaches typically run on a connected laptop or server, not on the wearable hardware itself.

\textbf{Deep learning on IMU time-series} has demonstrated superior accuracy without manual feature engineering. Nho et al.\ \cite{nho2021cnn} trained a 1D CNN on the UCI MHEALTH dataset and SisFall combined, reporting sensitivity of 97.2\% and specificity of 98.1\% with a model that fits within 150\,KB. However, the model ran on a Raspberry Pi 3 serving as the wearable compute unit, not on a resource-constrained microcontroller. Wang et al.\ \cite{wang2020lstm} applied an LSTM architecture to continuous IMU streams, achieving comparable accuracy but with inference latency unsuitable for real-time detection on embedded hardware.

\textbf{TFLite Micro on microcontrollers} has enabled CNN deployment on devices with fewer than 320\,KB of RAM. David et al.\ \cite{david2021tensorflow} demonstrated keyword spotting and gesture recognition using TFLite Micro on Cortex-M4 class hardware, establishing the feasibility of INT8 quantized CNN inference in this memory envelope. Kumar et al.\ \cite{kumar2022fall} deployed a threshold-gated classifier on an Arduino Nano~33~BLE Sense, but did not combine on-device TFLite Micro CNN inference with a structured two-layer gating architecture or provide any gateway explainability.

\textbf{Explainability in health AI} has been identified as a critical gap for clinical adoption. Lundberg and Lee \cite{lundberg2017shap} introduced SHAP (SHapley Additive exPlanations), a unified framework for model interpretability based on cooperative game theory, which provides consistent attribution of each input feature's contribution to a model output. While SHAP has been applied to ECG classification \cite{yao2021interpretation} and sepsis prediction in clinical ML, no published wearable fall detection system applies SHAP attribution at the per-event level and surfaces the result to caregivers in a clinical dashboard.

\textbf{Gap summary:} The existing literature leaves the following cluster of gaps unaddressed simultaneously: (a) on-device TFLite Micro CNN inference in a two-layer gated architecture on a sub-\$5 microcontroller; (b) SHAP-based per-event explainability surfaced to caregivers; (c) adaptive sensitivity profiles without firmware reflash; and (d) a dataset collected from Nepali subjects on locally procurable hardware. SPARK addresses all four gaps in a single open-source system.

\begingroup
\small
\setlength{\tabcolsep}{4pt}
\begin{longtable}{|p{0.7in}|p{1.4in}|p{1.5in}|p{1.5in}|}
\caption{Literature Review Matrix} \label{tab:litreview} \\
\hline
\textbf{Authors} & \textbf{Paper Title (Year)} & \textbf{Key Contribution} & \textbf{Motivation / Adopted Concepts} \\
\hline
\endfirsthead

\multicolumn{4}{c}{\tablename\ \thetable\ -- \textit{continued from previous page}} \\
\hline
\textbf{Authors} & \textbf{Paper Title (Year)} & \textbf{Key Contribution} & \textbf{Motivation / Adopted Concepts} \\
\hline
\endhead

\hline
\multicolumn{4}{r}{\textit{continued on next page}} \\
\endfoot

\hline
\endlastfoot

Kangas et al.\ \cite{kangas2008comparison} &
Comparison of Low-Complexity Fall Detection Algorithms for Body-Attached Accelerometers (2008) &
Compares threshold algorithms on hip/wrist IMU data; 73--95\% sensitivity depending on placement, false-positive rate $>$30\% on vigorous ADLs. &
Motivated Layer 1 threshold gating in SPARK's two-layer architecture, with awareness that threshold alone is insufficient. \\
\hline

Bourke \& Lyons \cite{bourke2008threshold} &
Evaluation of a Threshold-Based Tri-Axial Accelerometer Fall Detection Algorithm (2008) &
Dual-axis threshold method reaching 96\% sensitivity; specificity poor on sit-to-stand transitions. &
Confirmed the need for a learned classifier stage beyond threshold gating alone. \\
\hline

Sucerquia et al.\ \cite{sucerquia2017sisfall} &
SisFall: A Fall and Movement Dataset (2017) &
Public benchmark: 15 fall types, 19 ADL types, 38 subjects. &
Adopted as SPARK's primary training dataset, paired with self-collected Nepal-context data. \\
\hline

Nho, Lim \& Kwon \cite{nho2021cnn} &
Deep Convolutional Neural Network for Fall Detection Using Inertial Sensor Data from Wearable Devices (2021) &
1D CNN on UCI MHEALTH + SisFall; 97.2\% sensitivity, 98.1\% specificity, model fits in 150\,KB, but ran on Raspberry Pi 3. &
Motivated SPARK's on-device TFLite Micro target instead of an SBC-class compute unit. \\
\hline

Wang et al.\ \cite{wang2020lstm} &
LSTM-Based Fall Detection System Using Wearable Sensors (2020) &
LSTM on continuous IMU streams; accuracy comparable to CNN but with latency unsuitable for real-time embedded use. &
Informed SPARK's choice of CNN over LSTM for the MCU inference stage. \\
\hline

David et al.\ \cite{david2021tensorflow} &
TensorFlow Lite Micro: Embedded Machine Learning for TinyML Systems (2021) &
Demonstrates INT8-quantized CNN inference under 320\,KB RAM on Cortex-M4-class hardware (keyword spotting, gesture recognition). &
Directly adopted as SPARK's inference runtime (ESP32-S3 + TFLite Micro CNN). \\
\hline

Kumar, Sharma \& Singh \cite{kumar2022fall} &
Microcontroller-Based Fall Detection for Elderly Monitoring Using a Threshold-Gated Classifier on Arduino Nano 33 BLE Sense (2022) &
Threshold-gated classifier deployed and evaluated entirely on-device, but without TFLite Micro CNN inference or explainability. &
Motivated closing this exact gap: on-device CNN inference plus gating plus explainability, together. \\
\hline

Lundberg \& Lee \cite{lundberg2017shap} &
A Unified Approach to Interpreting Model Predictions (2017) &
Introduces SHAP, a cooperative game-theoretic framework for consistent per-feature attribution. &
Adopted as SPARK's explainability layer for per-event attribution surfaced to caregivers. \\
\hline

Yao et al.\ \cite{yao2021interpretation} &
Interpretation of Electrocardiogram Heartbeat by CNN and GRU (2021) &
Applies SHAP to a clinical CNN+GRU ECG classifier; attribution highlighted physiologically plausible regions. &
Confirms SHAP generalizes to real clinical deep-learning models; no prior fall-detection work applies it at the per-event level. \\
\hline

\end{longtable}
\endgroup